\subsection{Cohomology of Absolutely Semisimple Algebras}

PROPOSITION 6. --- Let K be a commutative ring and A a K-algebra. Let \(e = \sum_{i=1}^{r} a_i \otimes a_i'\) be an element of \(B = A \otimes_K A^o\) satisfying the conditions in item (ii) of Proposition 5 of VIII, p.~236. For any integer \(n \geqslant 1\) and any element f of \(C^n(A, P)\) , we denote by \(\gamma^n(f)\) the element of \(C^{n-1}(A, P)\) defined by the formula

\[
\gamma^ {n} (f) \left(x _ {1}, \dots , x _ {n - 1}\right) = \sum_ {i = 1} ^ {r} a _ {i} f \left(a _ {i} ^ {\prime}, x _ {1} \dots , x _ {n - 1}\right).\tag{15}
\]

We then have

\[
\partial^ {n - 1} (\gamma^ {n} (f)) + \gamma^ {n + 1} (\partial^ {n} (f)) = f\tag{16}
\]

for every integer \(n \geqslant 1\) and every \(f \in \mathrm{C}^n(\mathrm{A},\mathrm{P})\) .

*Remark 1. --- The morphisms \(\partial^n: \mathrm{C}^n(\mathrm{A},\mathrm{P}) \to \mathrm{C}^{n+1}(\mathrm{A},\mathrm{P})\) define a complex \((\mathrm{C}(\mathrm{A},\mathrm{P}),\partial)\) of K-modules (X, §2, n° 1, p.~24). The mapping \(\gamma_n\) therefore defines a homotopy from this complex to itself linking 0 to \(\mathrm{Id}_{\mathrm{C}(\mathrm{A},\mathrm{P})}\) (X, §2, n° 4, p.~32, définition 4). *

We keep the notation of No.6. For every integer \(n \geqslant 0\) , we define a mapping \(h_n: \mathrm{B}_n \to \mathrm{B}_{n+1}\) by the formula

\[
h _ {n} (x) = d _ {n + 2} ^ {1} (e \otimes x) = \sum_ {i = 1} ^ {r} a _ {i} \otimes a _ {i} ^ {\prime} x.
\]

It is a homomorphism of \((A, A)\) -bimodules (formula (3)).

Lemma 5. --- We have the relation

\[
d _ {n + 1} \circ h _ {n} + h _ {n - 1} \circ d _ {n} = 1 _ {\mathrm{B} _ {n}}\tag{17}
\]

for every \(n \geqslant 1\) .

Let \(x\in \mathrm{B}_n\) ; we have

\[
\begin{array}{c} (d _ {n + 1} \circ h _ {n}) (x) = (d _ {n + 1} \circ d _ {n + 2} ^ {1}) (e \otimes x) \\ = (d _ {n + 1} ^ {0} \circ d _ {n + 2} ^ {1}) (e \otimes x) - \sum_ {i = 2} ^ {n + 2} (- 1) ^ {i} (d _ {n + 1} ^ {i - 1} \circ d _ {n + 2} ^ {1}) (e \otimes x), \end{array}
\]

so, by formula (6),

\[
(d _ {n + 1} \circ h _ {n}) (x) = (d _ {n + 1} ^ {0} \circ d _ {n + 2} ^ {0}) (e \otimes x) - \sum_ {i = 2} ^ {n + 2} (- 1) ^ {i} (d _ {n + 1} ^ {1} \circ d _ {n + 2} ^ {i}) (e \otimes x).
\]

But we have

\[
(d _ {n + 1} ^ {0} \circ d _ {n + 2} ^ {0}) (e \otimes x) = \varepsilon (e) x = x
\]

by property (ii) of Proposition 5 of VIII, p.~236 and, for \(i \geqslant 2\) ,

\[
d _ {n + 2} ^ {i} (e \otimes x) = e \otimes d _ {n} ^ {i - 2} (x),
\]

which gives

\[
(d _ {n + 1} \circ h _ {n}) (x) = x - d _ {n + 1} ^ {1} (e \otimes d _ {n} (x)) = x - h _ {n - 1} \circ d _ {n} (x)
\]

and therefore formula (17).

We can finish the proof of Proposition 6 using Lemma 5. Let n be an integer \(\geqslant 1\) and f an element of \(C^{n}(A,P)\) . By construction, we have

\[
\alpha^ {n - 1} (\gamma^ {n} (f)) = \alpha^ {n} (f) \circ h _ {n - 1}\tag{18}
\]

and, consequently, by formulas (9) and (18),

\[
\begin{array}{r l} \alpha^ {n} (\partial^ {n - 1} (\gamma^ {n} (f)) + \gamma^ {n + 1} (\partial^ {n} (f))) & = \alpha^ {n - 1} (\gamma^ {n} (f)) \circ d _ {n} + \alpha^ {n + 1} (\partial^ {n} (f)) \circ h _ {n} \\ & = \alpha^ {n} (f) \circ h _ {n - 1} \circ d _ {n} + \alpha^ {n} (f) \circ d _ {n + 1} \circ h _ {n} \\ & = \alpha^ {n} (f), \end{array}
\]

where the last equality follows from (17). Since \(\alpha^{n}\) is bijective, the proposition follows.

THEOREM 3. --- Let K be a commutative ring, A a K-algebra, and P an (A,A)-bimodule. Suppose that the \((\mathrm{A}\otimes_{\mathrm{K}}\mathrm{A}^{\circ})\) -module A is projective. We then have \(\mathrm{H}^{n}(\mathrm{A},\mathrm{P})=0\) for every integer \(n\geqslant1\) .

We must prove that for every integer \(n \geqslant 1\) , every element f of \(C^{n}(A, P)\) with \(\partial^{n}(f) = 0\) is of the form \(\partial^{n-1}(g)\) for an element g of \(C^{n-1}(A, P)\) . By Proposition 5, this is an immediate consequence of Proposition 6.

COROLLARY. --- Every K-derivation from A to P is inner.

This is a translation of the equality \(H^{1}(A,P)=0\) .

Remarks. --- 2) The assumptions of Theorem 3 are, in particular, satisfied when K is a field and A an absolutely semisimple K-algebra (VIII, p.~238, Theorem 2).

*3) Suppose that the K-module A is projective. Theorem 3 can also be proved as follows. The complex \((\bigoplus_{n\geqslant 0}\mathrm{B}_n,d)\) and the homomorphism \(\varepsilon :\mathrm{B}_0\to \mathrm{A}\) define a projective resolution of the B-module A; therefore, for every \(n\geqslant 0\) , the K-module \(\mathrm{H}^n (\mathrm{A},\mathrm{P})\) is isomorphic to \(\operatorname {Ext}_{\mathrm{B}}^{n}(\mathrm{A},\mathrm{P})\) (X, §6, n° 1, p.~100, théorème 1). If the B-module A is projective, then the K-modules \(\operatorname {Ext}_{\mathrm{B}}^{n}(\mathrm{A},\mathrm{P})\) are zero for \(n\geqslant 1\) (X, §5, n° 3, p.~88, corollaire de la proposition 5), which implies that \(\mathrm{H}^n (\mathrm{A},\mathrm{P})\) is zero. Conversely, if \(\mathrm{H}^1 (\mathrm{A},\mathrm{P})\) is zero for every (A, A)-bimodule P, then the B-module A is projective (X, §5, n° 5, p.~93, proposition 10).*

